\chapter{Background and Related Work}
\label{ch:background}

The decision to use reinforcement learning changed the problem rather than removing it. Instead
of writing a complete controller by hand, I now had to build an environment in which useful
control could be learned. Four ideas shaped that environment: the physical demands of powered
descent, the reinforcement-learning formulation, the behaviour of PPO, and the distinction
between progress during training and performance during evaluation.

\section{Reusable-booster control}

A Falcon-class booster has nine first-stage engines, four deployable landing legs and four grid
fins, and burns liquid oxygen and rocket-grade kerosene \cite{spacex2025falcon}. These details
provided a useful scale and suggested which components the simulator would eventually need.
They did not make the implemented vehicle a certified Falcon~9 model or digital twin: public
dimensions and propulsion data establish plausible proportions, while the controller learns
only from the physical model defined in Chapter~\ref{ch:physics}.

The landing problem is difficult because translation and rotation cannot be solved separately.
Tilting the vehicle redirects thrust to correct a horizontal error, but simultaneously reduces
the vertical component needed to slow the descent. The amount of available control then changes
with throttle limits, actuator delay, fuel consumption, atmosphere and contact with the ground.
What appears visually to be one smooth manoeuvre is therefore a sequence of tightly coupled
decisions made under a shrinking margin for error.

Powered-descent guidance is often formulated as an optimal-control problem. Convexification can
transform relevant constraints into a form suitable for onboard optimization
\cite{acikmese2007convex}. Neural methods can instead approximate a control law after observing
many solved or sampled trajectories \cite{sanchez2018dnncontrol}, while deep reinforcement
learning has been applied directly to six-degree-of-freedom planetary landing with continuous
actions and randomized initial conditions \cite{gaudet2020landing}. This thesis shares the
control objective of that work, but asks a different engineering question: how can an
interactive Unity environment remain configurable and reproducible while the vehicle, task and
reward are still evolving?

\section{Reinforcement learning model}

An episodic Markov decision process is described by states $s_t$, actions $a_t$, transition
dynamics, rewards $r_t$ and a discount factor $\gamma$. A policy
$\pi_\theta(a_t\mid s_t)$ with parameters $\theta$ selects actions so as to maximize the expected
discounted return
\begin{equation}
  J(\theta)=\mathbb{E}_{\pi_\theta}\!\left[\sum_{t=0}^{T-1}\gamma^t r_t\right].
  \label{eq:return}
\end{equation}
The simulator does not provide the policy with a perfect physical state. It constructs an
observation vector $o_t$ from normalized relative kinematics, attitude and actuator state, so
the learned controller is more precisely written $\pi_\theta(a_t\mid o_t)$.

This formulation makes the reward function especially important. A human description such as
``land softly near the centre'' is not available to the optimizer; only the numerical reward
and termination rules are. A collection of individually sensible terms can therefore describe
an unintended strategy more profitably than the intended one. Potential-based shaping can
preserve an optimal policy under specific conditions \cite{ng1999policy}, but this project also
uses bounded physical costs, event rewards and terminal outcomes. Those terms have to be audited
empirically. The early policy that hovered just above the ground illustrates why: a small reward
paid repeatedly eventually outweighed the one-off value of actually landing.

\section{Proximal Policy Optimization}

PPO was selected because Unity ML-Agents supports it directly and because the rocket primarily
uses continuous control. PPO alternates between collecting trajectories under an existing
policy and improving that policy through a clipped surrogate objective
\cite{schulman2017ppo}. With probability ratio
\begin{equation}
  r_t(\theta)=\frac{\pi_\theta(a_t\mid o_t)}
  {\pi_{\theta_{\mathrm{old}}}(a_t\mid o_t)},
\end{equation}
the clipped objective is
\begin{equation}
  L^{\mathrm{CLIP}}(\theta)=
  \mathbb{E}_t\left[\min\left(r_t(\theta)\hat A_t,
  \operatorname{clip}(r_t(\theta),1-\epsilon,1+\epsilon)\hat A_t\right)\right],
  \label{eq:ppo}
\end{equation}
where $\hat A_t$ estimates the advantage of an action. Clipping limits the incentive for one
update to move too far from the policy that produced the data.

That stability is useful, but it does not protect the experiment from a missing observation, an
incorrect force or a reward loophole. PPO will optimize the task it is given, including any
mistake in that task. The experiments therefore use it as a consistent way to validate the
simulator rather than as a claim that PPO is the best algorithm for rocket landing. Unity
ML-Agents exchanges trajectories with the Python trainer, and the resulting multilayer
perceptron is exported to ONNX for inference inside Unity. No camera input is used; all policy
inputs are numerical observations derived from the simulated state.

\section{Curricula and evaluation}

The complete landing and tracking tasks were too broad to present to an untrained policy at
once. The simulator therefore represents difficulty with a scalar $d\in[0,1]$ and interpolates
between initial and final parameter ranges,
\begin{equation}
  x(d)=(1-d)x_{\mathrm{initial}}+d x_{\mathrm{full}}.
  \label{eq:curriculum-interpolation}
\end{equation}
As performance improves, the curriculum increases spawn distances, initial velocities or the
strictness of success thresholds. The sampled difficulty is frozen for the duration of an
episode, or for one target in the tracking task, so that the rules cannot change midway through
an attempt. Easier replay episodes remain in the distribution to reduce forgetting.

Curriculum progress answers whether the trainer is making the task harder; it does not answer
how reliable the exported policy is. Progress depends on recent training samples, adaptive
counters and stochastic actions. This became important in L11, where exact-full-difficulty
training success was much higher than deterministic evaluation success. For that reason,
evaluation fixes the seed, tests explicit difficulty bands and records every episode separately.

Binary success is reported with a Wilson score interval. For $n$ trials, observed proportion
$\hat p$ and normal quantile $z=1.96$, the interval is
\begin{equation}
\frac{\hat p+z^2/(2n)\ \pm\ z\sqrt{\hat p(1-\hat p)/n+z^2/(4n^2)}}
     {1+z^2/n}.
\label{eq:wilson}
\end{equation}
It remains well behaved near zero and one, where a symmetric normal approximation is misleading.

Deep-RL results can also vary across training seeds and implementation details
\cite{henderson2018matters,agarwal2021precipice}. Because the available computation and schedule
allowed only one principal training seed per task, the numerical results in this thesis validate
the simulator workflow and describe these particular policies. They are not a statistically
general comparison between algorithms.
